% ============================================================
%  Mission 50' — Nombres complexes
%  Carnet de bord — Équipe PAPIER (sans ordinateur)
%  Compilation : pdflatex (TeXstudio : F5, ouCompilation rapide)
% ============================================================
\documentclass[11pt,a4paper]{article}

\usepackage[utf8]{inputenc}
\usepackage[T1]{fontenc}
\usepackage[french]{babel}
\usepackage{lmodern}
\usepackage{amsmath,amssymb}
\usepackage[margin=1.8cm]{geometry}
\usepackage{tabularx}
\usepackage{enumitem}
\usepackage{tikz}
\usepackage{xcolor}
\usetikzlibrary{arrows.meta}

% ---------- Réglages ----------
\setlength{\parindent}{0pt}
\setlength{\parskip}{4pt}
\renewcommand{\arraystretch}{1.4}

% Grille vierge : origine O au centre, 1 carreau = 1 unite
% #1 = borne en x, #2 = borne en y
\newcommand{\grille}[3][1 carreau = 1 unité]{%
\begin{center}
\begin{tikzpicture}[scale=0.8,>=Stealth]
  \draw[very thin,gray!40] (-#2,-#3) grid (#2,#3);
  \draw[->,thick] (-#2,0) -- (#2+0.5,0) node[right]{réels};
  \draw[->,thick] (0,-#3) -- (0,#3+0.5) node[above]{imaginaires};
  \fill (0,0) circle (2.6pt);
  \node[below left] at (0,0) {$O$};
  \node[below] at (current bounding box.south) {\footnotesize Origine $O$ au centre --- #1};
\end{tikzpicture}
\end{center}}

% Ligne de réponse à compléter
\newcommand{\pointilles}{\makebox[3.5cm]{\dotfill}}

\begin{document}

% ============================================================
%  PAGE 1 — CARNET DE BORD
% ============================================================
\begin{center}
  {\LARGE\bfseries Carnet de bord --- équipe PAPIER}\\[2pt]
  {\large (sans ordinateur)}
\end{center}

Imprimez cette section pour l'équipe papier. Le bouton régénère toutes les
valeurs : cliquez avant chaque impression pour donner une série différente à
chaque binôme.

\textbf{Graine de cette série : 670421} --- notez-la sur la fiche : l'enseignant
peut rejouer exactement ces mêmes valeurs.

\section*{Équipe PAPIER --- 10 élèves, 5 binômes, sans ordinateur}

\begin{itemize}
  \item L'enseignant imprime le carnet papier de ce document (section
        « Carnet ») : fiches A à F + feuille rallye. Cliquez d'abord sur
        \emph{régénérer} pour changer toutes les valeurs, et imprimez une série
        différente par binôme si vous voulez.
  \item \textbf{Phase 1} --- fiche A (l'histoire en 5 dates) + fiche C
        (chasse aux applications dans le quotidien).
  \item \textbf{Phase 2} --- fiche B (le duel de calcul) + fiche D
        (la chasse au trésor dans le plan).
  \item \textbf{Phase 3} --- feuille rallye : chaque binôme fabrique 2 questions
        (une algébrique, une trigonométrique) pour l'équipe Écran, avec sa
        grille de correction.
  \item \textbf{Rôles par binôme :} un calculateur, un contrôleur (vérifie,
        trace, mesure), on échange à la phase 2.
\end{itemize}

\subsection*{Ce que votre binôme rend (tout est dans ce carnet)}

\begin{itemize}
  \item Le carnet complété : fiches A, B, C, D (et E pour les rapides) ;
  \item La feuille rallye : vos 2 questions fabriquées avec leur grille de
        correction, validées par l'enseignant ;
  \item La fiche F : vos 3 cartes de transmission remplies, pour la passerelle ;
  \item La feuille 4 du canevas distribué (fiche bilan individuelle) : chacun
        la remplit et la conserve pour le prochain cours.
\end{itemize}

\subsection*{Pendant ce temps, l'équipe Écran (ordinateur) découvre\dots}

\begin{itemize}
  \item les puissances de $i$ tournent en boucle : $i$, $-1$, $-i$, $1$, puis
        cela recommence (période 4) ;
  \item multiplier par $i$ semble faire tourner le point d'un quart de tour sans
        changer sa distance à l'origine ;
  \item la formule $(a + ib)(c + id) = (ac - bd) + (ad + bc)\,i$ résiste à
        tous les essais aléatoires.
\end{itemize}

Autant d'indices pour orienter vos tracés : à vous de les \textbf{confirmer par
la géométrie} (fiches B et E).

\newpage
% ============================================================
%  FICHE A — L'histoire en 5 dates
% ============================================================
\section*{FICHE A --- L'histoire en 5 dates}

Binôme : \pointilles\ --- Lecture 8~min, questions 5~min.

\textbf{1545, Cardan.} Les mathématiciens italiens s'affrontent en défis
publics pour résoudre les équations du 3\up{e} degré. La formule de Cardan
marche\dots{} mais certains cas forcent à écrire $\sqrt{-121}$ :
l'« impossible ». On préfère ne pas en parler.

\textbf{1572, Bombelli.} Il a une idée « sauvage » : traiter $\sqrt{-1}$
comme un nombre ordinaire, avec les règles habituelles. Les monstres
s'éliminent entre eux\dots{} et donnent les solutions exactes, bien réelles,
de l'équation. On garde le calcul sans le comprendre.

\textbf{1637, Descartes.} Il les baptise « imaginaires », par dédain : ils
n'existeraient que dans l'imagination. Le nom est resté, le dédain non.

\textbf{1806, Argand.} Le déclic : représenter $a + ib$ par un point dans un
plan. Ajouter devient additionner des flèches ; multiplier devient tourner et
agrandir. Le monstre a un visage.

\textbf{1831, Gauss.} Il donne la théorie rigoureuse, un nouveau nom ---
« complexes » --- et un théorème : toute équation polynomiale a une solution
dans $\mathbb{C}$. Statut officiel. Deux siècles plus tard, c'est l'outil
standard de l'électricien et de l'ingénieur.

\textbf{Questions} (répondez sur cette fiche, entourez ou écrivez) :
\begin{enumerate}
  \item Quelle équation, toute simple, les nombres réels ne savent-ils pas
        résoudre ?
  \item Qui a le premier osé calculer avec $\sqrt{-1}$, et en quelle année ?
  \item Qui a eu l'idée de représenter un complexe par un point d'un plan ?
  \item Que veut dire le mot « complexe » choisi par Gauss ?
  \item Retenez \textsc{une} phrase d'histoire à raconter au reste de la
        classe.
\end{enumerate}

\newpage
% ============================================================
%  FICHE B — Le duel de calcul
% ============================================================
\section*{FICHE B --- Le duel de calcul}

Binôme : \pointilles\ --- \emph{Calculatrice autorisée. 12~min.}

\textbf{Règle :} pour chaque ligne du tableau, calculateur et contrôleur se
partagent les colonnes, puis comparent. Le but : finir la ligne la plus juste.
Envie de lignes bonus ? Créez vos propres $z$ en lançant deux dés : un dé rouge
donne $a$, un dé blanc donne $b$ ($6 = 0$, et un 1 relancé = nombre négatif).

\begin{center}
\begin{tabular}{|c|c|c|c|c|c|}
  \hline
  \textbf{Ligne} & $z_A$ & $z_B$ & $z_A + z_B$ & $z_A \times z_B$
  & $|z_A|$ (arrondi 0,1) \\ \hline
  1 & $5 + 2i$ & $-5 + 3i$ & & & \\ \hline
  2 & $1 + i$  & $2 + i$   & & & \\ \hline
  3 & $4 - 4i$ & $-3 + 2i$& & & \\ \hline
  4 & $5 + 2i$ & $5 - 2i$ & & & \\ \hline
  5 & $2 - 4i$ & $3 + 4i$ & & & \\ \hline
  6 & $-1 + 4i$& $5 + 2i$ & & & \\ \hline
\end{tabular}
\end{center}

\textbf{À tracer :} placez les points $z_A$ et $z_B$ de la ligne 1 sur la
grille ci-dessous (1 carreau = 1 unité, l'origine est marquée).

\grille[1 carreau = 1 unité ; axe des réels $\rightarrow$, axe des imaginaires
$\uparrow$]{9}{7}

\newpage
% ============================================================
%  FICHE C — La chasse aux applications
% ============================================================
\section*{FICHE C --- La chasse aux applications}

Binôme : \pointilles\ --- Brainstorm 12~min. \emph{Sans ordinateur :} votre
matière grise, vos souvenirs, le lycée, la maison.

\textbf{Consigne :} listez un maximum d'objets, d'appareils ou de technologies
du quotidien qui, selon vous, utilisent des nombres complexes (courant
alternatif, ondes, son, image, moteur, capteur\dots). Pour chacun, notez en
deux mots ce qui « tourne » ou « oscille » dedans. Puis classez : votre
\textbf{top 3} est à défendre à l'oral devant l'équipe Écran à la phase 4.

\begin{center}
\renewcommand{\arraystretch}{2.2}
\begin{tabularx}{\textwidth}{|c|X|X|c|}
  \hline
  \textbf{\#} & \textbf{Objet / technologie}
  & \textbf{Ce qui oscille ou tourne dedans\dots} & \textbf{Vérifié ?} \\
  \hline
  1 & & & \\ \hline
  2 & & & \\ \hline
  3 & & & \\ \hline
  4 & & & \\ \hline
  5 & & & \\ \hline
  6 & & & \\ \hline
  7 & & & \\ \hline
  8 & & & \\ \hline
\end{tabularx}
\end{center}

\textbf{Notre top 3 à défendre :}\newline
1)~\pointilles\newline 2)~\pointilles\newline 3)~\pointilles

\textbf{Notre argument en une phrase :}

\pointilles\par
\pointilles

\newpage
% ============================================================
%  FICHE D — La chasse au trésor dans le plan complexe
% ============================================================
\section*{FICHE D --- La chasse au trésor dans le plan complexe}

Binôme : \pointilles\ --- 12~min. Départ à l'origine $O$. Le trésor est caché
au point d'arrivée.

\textbf{Consigne :} chaque étape est un déplacement : « $+(3 + 2i)$ »
signifie avancer de 3 vers l'est et de 2 vers le nord. Additionner deux
complexes, c'est additionner les déplacements, composante par composante.
Suivez le parcours, tracez-le, et répondez.

\begin{center}
\begin{tabular}{|c|c|c|}
  \hline
  \textbf{Étape} & \textbf{Déplacement} & \textbf{Position $a + ib$ après l'étape} \\
  \hline
  1 & $+(-3 - 3i)$ & \\ \hline
  2 & $+(2 + 3i)$  & \\ \hline
  3 & $+(-2 - 3i)$ & \\ \hline
  4 & $+(-3 - i)$  & \\ \hline
  5 & $+(1 + 2i)$  & \\ \hline
  6 & $+(2 - 3i)$  & \\ \hline
  7 & $+(3 + 2i)$  & \\ \hline
  8 & $+(2 + i)$   & \\ \hline
\end{tabular}
\end{center}

\begin{enumerate}
  \item Tracez le parcours sur la grille (départ $O$, 1 carreau = 1 unité).
  \item Quelles sont les coordonnées $a$ et $b$ du point du trésor ?
  \item Quelle est la distance à vol d'oiseau entre $O$ et le trésor (le
        module) ? Arrondi au dixième.
  \item À votre avis : quel angle fait la direction $OT$ avec l'axe des
        réels ? (estimez au rapporteur)
  \item \textbf{Bonus :} que se passerait-il si, au lieu d'ajouter la dernière
        étape, on multipliait la position par $i$ ? (indice : rotation\dots)
\end{enumerate}

\grille[1 carreau = 1 unité --- départ de la chasse]{9}{7}

\newpage
% ============================================================
%  FICHE E — La preuve du géomètre
% ============================================================
\section*{FICHE E --- La preuve du géomètre (défi, pour les rapides)}

À deux, avec la grille de la fiche B. 8~min. C'est votre « preuve visuelle »
pour la passerelle.

On sait (fiche B) que $i(a + ib) = ia + i^2 b = -b + ai$. Donc si le point
$M(a\,;\,b)$ porte l'affixe $z$, le point $M'(-b\,;\,a)$ porte l'affixe $iz$.
Prouvez que c'est un quart de tour :

\textbf{Étape 1 --- les modules.} Calculez $OM^2 = a^2 + b^2$ et
$OM'^2 = (-b)^2 + a^2$, comparez, puis concluez sur $|z|$ et $|iz|$.

\begin{center}
\renewcommand{\arraystretch}{2.2}
\begin{tabular}{|c|c|c|}
  \hline
  $OM^2$ & $OM'^2$ & Conclusion sur les modules \\ \hline
  & & \\ \hline
\end{tabular}
\end{center}

\textbf{Étape 2 --- l'angle.} Le produit scalaire des vecteurs
$\overrightarrow{OM}$ et $\overrightarrow{OM'}$ vaut $xx' + yy'$, soit ici
$a \times (-b) + b \times a$. Calculez-le. Que vaut-il ? Qu'en déduit-on pour
l'angle $\widehat{MOM'}$ ?

\begin{center}
\renewcommand{\arraystretch}{2.2}
\begin{tabular}{|c|c|}
  \hline
  $\overrightarrow{OM} \cdot \overrightarrow{OM'}$
  & Angle $\widehat{MOM'}$ déduit \\ \hline
  & \\ \hline
\end{tabular}
\end{center}

Conclusion à écrire : multiplier un nombre complexe par $i$, c'est lui faire
subir \pointilles\ (module \pointilles).

\newpage
% ============================================================
%  FICHE F — Cartes de transmission
% ============================================================
\section*{FICHE F --- Cartes de transmission (à garder pour la phase 4)}

Rédigez 3 cartes, une par ligne : \textcircled{1} une phrase courte (la
découverte) \,\textperiodcentered\, \textcircled{2} une preuve (valeur
mesurée, tracé ou citation historique) \,\textperiodcentered\,
\textcircled{3} à quoi ça sert. Vous les présenterez oralement
(= 1 minute par carte) à l'équipe Écran lors de la passerelle.

\begin{center}
\renewcommand{\arraystretch}{3}
\begin{tabularx}{\textwidth}{|c|X|}
  \hline
  \textbf{Carte} & \textcircled{1} Découverte \,\textperiodcentered\,
  \textcircled{2} Preuve \,\textperiodcentered\, \textcircled{3} Ça sert à
  quoi ? \\ \hline
  1 & \\[1.6cm] \hline
  2 & \\[1.6cm] \hline
  3 & \\[1.6cm] \hline
\end{tabularx}
\end{center}

\newpage
% ============================================================
%  FEUILLE RALLYE — Fabricateurs de questions (phase 3)
% ============================================================
\begin{center}
  {\Large\bfseries FEUILLE RALLYE}\\[2pt]
  {\large Fabricateurs de questions (phase 3)}
\end{center}

Chaque binôme fabrique 2 questions pour l'équipe Écran et les fait valider par
l'enseignant. 10~min.

\textbf{Question 1 (forme algébrique).} Inventez un calcul du type « calculez
$z_1 + z_2$, $z_1 \times z_2$ ou $|z|$ » avec vos propres nombres (évitez les
valeurs trop faciles). Écrivez l'énoncé, puis la réponse détaillée en dessous.

\begin{center}
\renewcommand{\arraystretch}{3.5}
\begin{tabularx}{\textwidth}{|X|X|}
  \hline
  \textbf{Énoncé (à dicter à l'équipe Écran)}
  & \textbf{Réponse détaillée (votre correction)} \\ \hline
  & \\[2.2cm] \hline
\end{tabularx}
\end{center}

\textbf{Question 2 (forme trigonométrique).} Choisissez un angle remarquable
($30^{\circ}$, $45^{\circ}$, $60^{\circ}$, $120^{\circ}$, $135^{\circ}$,
$150^{\circ}$) et un entier $r$ entre 2 et 10. Demandez par exemple la forme
algébrique de $z = r(\cos\theta + i\sin\theta)$, ou l'inverse.

\begin{center}
\renewcommand{\arraystretch}{3.5}
\begin{tabularx}{\textwidth}{|X|X|}
  \hline
  \textbf{Énoncé} & \textbf{Réponse détaillée (votre correction)} \\ \hline
  & \\[2.2cm] \hline
\end{tabularx}
\end{center}

Puis échangez : l'équipe Écran répond à vos questions, vous répondez aux
siennes. Le binôme qui obtient le meilleur total marque la manche.

\end{document}
